\section{Related Work}
The human cost of conducting SLRs is concentrated in manual retrieval, screening, and evidence structuring~\citep{page2021prisma,marshall2019toward}. Early automation targeted study identification through machine-learned classifiers such as the Cochrane RCT Classifier~\citep{thomas2021machine}, and active-learning tools for screening \citep{gates2018technologyAbstrackr,przybyla2018prioritising, chai2021research} and risk-of-bias assessment \citep{gates2018technologyRobotReviewer}.
%, van2021open

Recent work with LLMs shows that prompt templates can transfer screening logic across title, abstract and full-text stages without task-specific fine-tuning, achieving high sensitivity and specificity in multiple systematic reviews~\citep{ottoSRPromptPaper,homiar2025development}. However, performance remains sensitive to class imbalance, prompt formulation and evolving eligibility criteria~\citep{khraisha2024can,syriani2024screening}. For data extraction, LLMs perform well on constrained schemas but degrade on complex fields, with human-incorporated LLM workflows generally outperforming LLM-only approaches~\citep{gartlehner2024data,mahmoudi2025critical,lai2025language}.

Building on these stage-specific advances, recent work has shifted towards end-to-end SLR pipelines that couple retrieval, screening, extraction and synthesis under agentic orchestration \citep{scherbakov2025emergence}. Some systems reproduce and update Cochrane-style intervention reviews by coordinating screening and extraction, but relies on proprietary models and evaluate performance using LLM-as-a-judge with post-hoc ``corrected'' labels \citep{ottosrCao2025automation}. Others emphasise full-text processing with traceable provenance and expert validation interfaces \citep{parkinson2025metabeeai} but do not fully automate upstream search and initial screening. Consequently, existing systems remain either proprietary or not targeted to WHO-designated priority pathogens.
